% Chapter 2: Foundational Technologies
\chapter{Foundational Technologies}
\label{ch:foundations}

This chapter describes the models, position-encoding machinery, and measurement toolkit
that the rest of the thesis builds on: the three evaluated diffusion VSR methods
(Section~\ref{sec:found:vsr}), rotary position embeddings and their extrapolation
behaviour (Section~\ref{sec:found:rope}), the perceptual and temporal quality measures
from which LR-VCC's sub-metrics are constructed (Section~\ref{sec:found:toolkit}), and
the DOVE evaluation protocol under which all ground-truth comparisons are performed
(Section~\ref{sec:found:dove}).

\section{Diffusion-Based Video Super-Resolution and the Evaluated Baselines}
\label{sec:found:vsr}

% TODO(Ch2 pass): 1 paragraph on latent diffusion for restoration (SD prior, VAE latent
% space, tile-based inference for large frames) before the per-model subsections.

\subsection{MGLD-VSR}
MGLD-VSR \cite{he_mgld_2024} (ECCV 2024) applies motion-guided latent diffusion to
progressively refine super-resolved frames, using optical-flow cues to steer the latent
denoising trajectory toward temporal coherence. It is characterised by sharp,
detail-preserving output. In this thesis it also serves as the numerically validated
anchor of the evaluation environment: it reproduces the published DOVE benchmark numbers
on UDM10 exactly (PSNR 24.23, with SSIM and LPIPS matching to four decimal places).

\subsection{Upscale-A-Video}
Upscale-A-Video (UAV) \cite{zhou_upscale_2024} (CVPR 2024) is a text-conditioned temporal
diffusion model inspired by text-to-video generation. It produces smoother, perceptually
cleaner output at the cost of some high-frequency detail. UAV reproduces DOVE UDM10
results within a consistent $+1.33$~dB PSNR offset attributable to a degradation-pipeline
configuration difference rather than a setup defect; the offset is consistent across test
sets ($+1.68$~dB on SPMCS), confirming the evaluation environment.

\subsection{FlashVSR}
FlashVSR \cite{zhuang_flashvsr_2025} is a one-step streaming VSR model distilled from the
Wan2.1 video diffusion transformer \cite{wan_wan_2025}. Frames are processed in a
causal streaming fashion with a KV cache carrying context across chunk boundaries, which
makes the method applicable to arbitrarily long input in principle. Its position encoding
is a three-axis factorised RoPE over (time, height, width), with the temporal index
growing monotonically over the whole stream --- the property that motivates the
extrapolation study of Chapter~\ref{ch:rope}.
% TODO(Ch6 prep): expand from docs/notes/2026-07-02-flashvsr-rope-site.md — DiT block
% structure, block-sparse attention with adaptive topk, TC decoder, distillation recipe,
% trained temporal window (21 latent frames ~ 81 video frames), spatial trained extent
% (48 latent), stock 1024-row frequency tables.

\section{Rotary Position Embeddings and Length Extrapolation}
\label{sec:found:rope}

Rotary position embeddings \cite{su_roformer_2024} encode a token's position $m$ by
rotating each two-dimensional channel pair of its query and key vectors by an angle
$m\theta_i$, where $\theta_i$ is a per-pair frequency drawn from a geometric schedule.
The attention logit between positions $m$ and $n$ then depends on positions only through
the relative offset $(m-n)$, giving translation invariance in expectation while retaining
absolute-position information within each head.

Video diffusion transformers factorise RoPE over three axes: the channel dimension is
partitioned into a temporal group and two spatial groups (height, width), each with its
own position index and frequency table. In Wan2.1-derived models such as FlashVSR the
tables are precomputed to a fixed number of rows; positions beyond the table are
unreachable at inference time unless the table is extended.

When a model is evaluated at positions or relative distances outside its trained range,
quality can degrade. The language-modelling literature offers three main remedies:
\emph{linear position interpolation} (PI) \cite{chen_pi_2023} compresses positions by a
constant factor so relative distances remain in-range; \emph{NTK-aware scaling}
\cite{bloc97_ntk_2023} rescales the frequency schedule non-uniformly, preserving
high-frequency (local) resolution; and \emph{YaRN} \cite{peng_yarn_2023} combines
frequency-dependent interpolation with attention-temperature correction. All three were
developed for one-dimensional text and validated on perplexity; their applicability to
three-axis video RoPE under streaming inference is one of the questions
Chapter~\ref{ch:rope} answers empirically.

\section{Perceptual and Temporal Quality Measurement Toolkit}
\label{sec:found:toolkit}

Three families of measures used throughout the thesis are described here.

\paragraph{No-reference perceptual measures.} CLIP-IQA \cite{wang_clipiqa_2023} evaluates
per-frame image quality using CLIP \cite{radford_learning_2021} embeddings. Unlike LPIPS
or DINOv2, CLIP-IQA does not systematically penalise high-frequency detail in
diffusion-style SR outputs, making it a suitable appearance component for a no-reference
composite. MUSIQ \cite{ke_musiq_2021} is a multi-scale image-quality transformer; DOVER
\cite{wu_dover_2023} predicts video quality through disentangled aesthetic and technical
branches.

\paragraph{Temporal consistency measures.} tOF and tLP \cite{chu_tecogan_2020} measure
optical-flow warping error and LPIPS-based perceptual warping error between frame pairs.
Conventionally both are evaluated at adjacent-frame lag $k{=}1$; this thesis extends them
to lags $k \in \{1, 5, 10, 30, 60, 120\}$ using RAFT \cite{teed_raft_2020} optical flow
with forward--backward consistency masking applied uniformly across all $k$.

\paragraph{Identity measures.} VBench-2.0's Human\_Identity dimension tracks face
embedding stability using RetinaFace \cite{deng_retinaface_2020} detection and ArcFace
\cite{deng_arcface_2019} embeddings. For long videos this thesis uses a slow-fast
adaptation: a \emph{slow} component fused within 2-second clips and a \emph{fast}
component comparing first-frame embeddings across clips, so that cross-clip identity
drift is measured rather than averaged away.
% TODO(Ch2 pass): mention the six upstream VBench-2.0 patches and the fps-override
% correction (hard-coded 30fps output tag) — one paragraph, with pointer to Ch3.

\section{The DOVE Evaluation Protocol}
\label{sec:found:dove}

All ground-truth comparisons in this thesis follow the DOVE \cite{yang_dove_2024}
protocol: low-resolution inputs are produced by the RealBasicVSR
\cite{chan_realbasicvsr_2022} degradation pipeline (bicubic downscale by a factor of
four, sensor noise injection, codec compression), and full-reference metrics are computed
in the DOVE RGB convention. The DOVE test suites used here are UDM10 \cite{yi_pfnl_2019}
and SPMCS \cite{tao_spmcs_2017} with realistic degradation, and YouHQ40
\cite{zhou_upscale_2024} for the resolution-extrapolation ladder of
Chapter~\ref{ch:rope}. Validating against DOVE's published numbers anchors every
experiment chapter: conclusions are drawn only on top of a numerically verified
inference environment.
